\vspace{-1mm}
\section{Conclusion}
\label{sec:conclusion}
Ensuring trust in LLM-powered data processing tasks is critical for maintaining usability. We presented \sys, a framework for inferring verifiable minimal provenance—subsets of the input text that reproduce an equivalent answer to the one obtained from the full input. \sys introduces a set of strategies, each guaranteed to return a minimal verifiable provenance, including an adaptive strategy that combines the strengths of multiple strategies, to infer provenance more efficiently. \techreport{In scenarios where a single provenance is insufficient to establish confidence in the answer's correctness, \sys offers a strategy to identify $k$ distinct minimal provenances, enabling users to gain stronger trust.} Extensive experiments across five datasets demonstrate that \sys is guaranteed to return a provenance that can reproduce the answer at a low cost, achieving over 30\% higher accuracy than the best-performing baseline. 

%So far, \sys focuses on inferring provenance for a single data processing task or operator. In future work, it would be valuable to explore provenance across pipelines in LLM-powered data systems. Key challenges include tracing provenance and its impact across multiple operators, formulating provenance to capture not only the relevant data portions but also the sequence of operations leading to the final result, and developing efficient strategies for provenance maintenance given the high cost of LLM invocations. 


 